\section{Introduction}\label{sec:introduction}

Viterbo's conjecture is an isoperimetric inequality in symplectic geometry
\cite{Viterbo2000}. In dimension four it states that every convex body
$K\subset\mathbb R^4$ satisfies
\[
 \operatorname{sys}(K)\coloneqq\frac{c_{\rm EHZ}(K)^2}{2\operatorname{vol}_4(K)}\leq1,
\]
with equality for the ball. Here $c_{\rm EHZ}(K)$, the
Ekeland--Hofer--Zehnder capacity, is the least action of a closed
characteristic on $\partial K$; for polytopes, generalized characteristics
replace closed characteristics. We call $\operatorname{sys}(K)$ the systolic
ratio. Haim-Kislev and Ostrover disproved the conjecture
\cite{HaimKislevOstrover2024}. Their counterexample is a Lagrangian product:
for planar convex bodies $K_q,K_p$, write $K_q\times_L K_p$ for the body with
$K_q$ in the $(q_1,q_2)$-plane and $K_p$ in the $(p_1,p_2)$-plane. With $P$ a
regular pentagon, the HKO body $K_{\rm HKO}=P\times_L R(-\pi/2)P$ has
systolic ratio $(3+\sqrt5)/5\approx1.047$.

This thesis picks a toolbox and tests how much progress it makes on
questions around Viterbo's conjecture after this counterexample. The toolbox
consists of high-performance computation on the LICCA cluster,
proof by computation in exact arithmetic, standard data-science methods,
and AI agents. The agents ran the data-science experiments and did large
parts of the theory work, including most of the local-maximality proof in
Section~\ref{sec:hko-local-maximum}. The questions were therefore not fixed
in advance; they came out of the work.

Data science has already contributed to pure mathematics. In knot theory,
patterns spotted in large tables of invariants led to results that then
inspired theory. We go one step further: AI agents apply the data-science
methods, and AI agents also do much of the resulting theory.

One of the earliest questions was whether standard data science finds new
interesting bodies with systolic ratio above one. It does not. Random samples
of generic polytopes and of Lagrangian products of polygons, and local
searches started from generic polytopes, stay below one. Bodies above one do
occur near $K_{\rm HKO}$, but none of them is new or interesting.
What the data did produce is a strong inverse association between systolic
ratio and a symplectic area of the two-dimensional faces, which the thesis
then examines with exact formulas and deformations.

Other questions arose the same way. Randomly sampled perturbations of
$K_{\rm HKO}$ never increased its systolic ratio. This suggested that
$K_{\rm HKO}$ is a local maximum, and the first question became whether this
holds. We prove it among polytopes with ten facets, and we find a second,
unrelated ten-facet local maximum with systolic ratio one. Rotating the
pentagon factors against each other led to an exact profile and to a sharp
bound for independently linearly deformed regular pentagons. The capacity
computations behind these results use Haim-Kislev's finite formula
\cite{HK2017}, a reduction for Lagrangian products, and a second algorithm
based on the local geometry of Chaidez and Hutchings \cite{CH2021}. The
results therefore answer different questions and do not form a single
narrative; Sections~1.1--1.5 introduce them in turn.
